\documentclass[a4paper,12pt]{article}

\usepackage[utf8]{inputenc}
\usepackage[T1]{fontenc}
\usepackage[brazil]{babel}
\usepackage{geometry}
\usepackage{enumitem}
\usepackage{longtable}
\usepackage{listings}

\geometry{margin=2.5cm}

\title{\textbf{RELATÓRIO FINAL -- COMPILADOR PYTHON DIDÁTICO}}
\author{
Marcelo Gabriel Euclides Da Silva -- Matrícula 01605103 \\
Giorgio Rafael Ataide Guimarães -- Matrícula 01620989 \\
Luiz Gustavo Freire Barbosa -- Matrícula 01864307 \\
Matheus Wendell de Paula Souza -- Matrícula 01602509 \\
Guilherme Oliveira Ribeiro -- Matrícula: 01606196 \\
Pedro Henrique Alves Lessa -- 01638620 \\
Matheus Venâncio Da Silva -- 01628436 \\
Yuri Mascarenhas -- 01327348 \\
Felipe de Lima Passos -- 01647492 \\
Maria Eduarda Cardoso Ferreira -- 01724547
}
\date{Maceió -- AL \\ 2026}

\begin{document}

\maketitle

\begin{center}
\textbf{Disciplina:} Compiladores \\[0.2cm]
\textbf{Curso:} Ciência da Computação
\end{center}

\newpage

\section{Visão Geral do Projeto}

\subsection{Problema a ser resolvido}

O estudo de compiladores é um dos temas mais importantes da Ciência da Computação, pois permite compreender como linguagens de programação são interpretadas e executadas pelos computadores. Entretanto, muitos estudantes encontram dificuldades para visualizar na prática as etapas envolvidas nesse processo, uma vez que grande parte do conteúdo é apresentada de forma teórica.

Diante desse cenário, foi desenvolvido um compilador didático inspirado na linguagem Python, permitindo que o usuário escreva código-fonte, realize sua análise e visualize os resultados da execução em uma interface gráfica intuitiva.

\subsection{Público-alvo}

O sistema foi desenvolvido para estudantes da disciplina de Compiladores, professores da área de Linguagens Formais e Compiladores, além de desenvolvedores iniciantes interessados em compreender o funcionamento interno de uma linguagem de programação.

\subsection{Proposta de Valor}

A proposta do projeto consiste em fornecer uma ferramenta educacional capaz de demonstrar, de forma prática, as principais etapas de um compilador moderno. O sistema realiza análise léxica, análise sintática e interpretação de código, permitindo ao usuário visualizar erros, executar programas e acompanhar o estado das variáveis durante a execução.

\section{Requisitos do Sistema}

\subsection{Requisitos Funcionais (RF)}

\begin{itemize}
    \item RF01: Permitir ao usuário escrever código-fonte em um editor integrado.
    \item RF02: Realizar análise léxica do código inserido.
    \item RF03: Realizar análise sintática utilizando uma gramática inspirada na linguagem Python.
    \item RF04: Executar programas válidos por meio de um interpretador interno.
    \item RF05: Exibir mensagens de erro contendo linha e coluna da ocorrência.
    \item RF06: Permitir a utilização de estruturas condicionais.
    \item RF07: Permitir a utilização de estruturas de repetição.
    \item RF08: Exibir as variáveis armazenadas em memória após a execução.
\end{itemize}

\subsection{Requisitos Não Funcionais (RNF)}

\begin{itemize}
    \item RNF01: O sistema deve ser desenvolvido utilizando JavaScript.
    \item RNF02: A interface gráfica deve ser desenvolvida utilizando React.
    \item RNF03: A aplicação deve possuir interface responsiva.
    \item RNF04: O sistema deve impedir execuções infinitas por meio de controle de passos.
    \item RNF05: A aplicação deve fornecer feedback visual sobre o resultado da compilação.
\end{itemize}

\section{Regras de Negócio}

\begin{itemize}
    \item Variáveis devem existir antes de serem utilizadas.
    \item Operações matemáticas devem ocorrer apenas entre valores numéricos compatíveis.
    \item Concatenação de textos deve ocorrer apenas entre valores textuais.
    \item O comando \texttt{range()} deve receber apenas valores inteiros.
    \item Loops infinitos devem ser interrompidos automaticamente quando o limite de execução for atingido.
    \item Comparações devem ocorrer apenas entre tipos compatíveis.
\end{itemize}

\section{Arquitetura do Sistema}

O sistema foi desenvolvido seguindo uma arquitetura modular, composta pelas etapas de análise léxica, análise sintática e interpretação de código.

\subsection{Fluxo de execução}

A interface desenvolvida em React é responsável pela entrada do código-fonte e pela exibição dos resultados. O módulo \texttt{Lexer} realiza a análise léxica, o \texttt{Parser} constrói a árvore sintática e o \texttt{Interpreter} executa as instruções geradas.

\section{Análise Léxica}

A análise léxica é a primeira etapa do processo de compilação. Seu objetivo é converter o código-fonte em uma sequência de tokens compreensíveis pelo compilador.

Entre os tokens suportados estão:

\begin{center}
\begin{tabular}{|l|l|}
\hline
\textbf{Código} & \textbf{Token} \\
\hline
if & IF \\
else & ELSE \\
while & WHILE \\
print & PRINT \\
true & BOOLEAN \\
\hline
\end{tabular}
\end{center}

\subsection{Exemplo}

Código-fonte:

\begin{lstlisting}
x = 10
\end{lstlisting}

Tokens gerados:

\begin{lstlisting}
IDENTIFIER EQUAL NUMBER
\end{lstlisting}

Além disso, o analisador léxico reconhece operadores matemáticos, operadores lógicos, operadores relacionais, números inteiros, números decimais, strings, caracteres e estruturas de indentação inspiradas na linguagem Python.

\end{document}